\chapter{Reranker}

\section{Identification}

\begin{center}
\begin{tabular}{ll}
\toprule
Modèle         & \texttt{BAAI/bge-reranker-v2-m3} (CrossEncoder) \\
Statut scoring & \textbf{Implémenté} \\
Fichier        & \texttt{DataPipelines/Reranker.py} \\
\bottomrule
\end{tabular}
\end{center}

\section{Fonctionnalité}

Le reranker reçoit les 25 chunks candidats des embeddings et les reclasse par
pertinence à la question en traitant chaque paire (question, chunk) ensemble.
Il filtre les chunks sous le seuil \texttt{minimum\_acceptable\_score = 0.05}
selon la méthode de sélection configurée (\texttt{smart}, \texttt{kmeans},
\texttt{top}, \texttt{top\_kmeans}).

\section{Fonction de scoring de confiance}

\noindent\textit{Implémentation dans \texttt{DataPipelines/Reranker.py}.}

\medskip
Le cross-encoder produit un logit $l_i$ par paire (question, chunk), converti
en score $\tilde{s}_i = \sigma(l_i) \in (0,1)$.

Le score de confiance actuel est une moyenne pondérée entre la moyenne des
chunks retenus et le plus faible d'entre eux :

\begin{equation}
  s_{\text{reranker}}
  = 0.7 \cdot \bar{s} + 0.3 \cdot \min_{i \in \mathcal{I}} \tilde{s}_i
  \label{eq:reranker}
\end{equation}

\noindent\textbf{Cette formule est temporaire.} La pondération 0.7/0.3 n'a
pas été validée empiriquement. Elle ne tient pas compte de la taille du corpus :
avec peu de documents disponibles, le système peut retenir des chunks
faiblement pertinents faute de mieux, ce que le score ne reflète pas.

\textbf{Cas particuliers.}
\begin{itemize}
  \item $|\mathcal{I}| = 0$ : $s = 0.0$ (aucun candidat retenu).
  \item $|\mathcal{I}| = 1$, fast path : $s = 0.5$ (non mesuré par convention).
\end{itemize}

\paragraph{Seuils.}

\begin{center}
\begin{tabular}{lll}
\toprule
$s \geq 0.70$ & Bons candidats & Réponse LLM probablement pertinente \\
$0.40 \leq s < 0.70$ & Pertinence moyenne & Vérification recommandée \\
$s < 0.40$ & Peu pertinent & Avertissement \\
\bottomrule
\end{tabular}
\end{center}

\noindent Ces seuils sont provisoires et non calibrés.

\textbf{Métadonnées loguées :} \texttt{num\_selected}, \texttt{mean\_score},
\texttt{min\_score}, \texttt{max\_score}, \texttt{per\_index\_scores}.


\noindent\textit{Tests de régression automatisés :
\texttt{Tests/test\_reranker\_regression.py}.}
